\documentclass[conference]{IEEEtran}
\IEEEoverridecommandlockouts
\usepackage{cite}
\usepackage{amsmath,amssymb,amsfonts}
\usepackage{booktabs}
\usepackage{multirow}
\usepackage{xcolor}
\usepackage{colortbl}
\usepackage{graphicx}
\usepackage{balance}
\usepackage[hidelinks]{hyperref}
\usepackage{enumitem}

% ---- collaboration macros ----------------------------------------------
% \wj{...}  : hypothesis/standpoint comment, rendered blue, greppable, strippable for camera-ready
% \ph{...}  : pre-registered placeholder awaiting data; every \ph is bound to the estimand block (Sec. IV-G)
\newcommand{\wj}[1]{\textcolor{blue}{\textbf{[WJ:} \textit{#1}\textbf{]}}}
\newcommand{\ph}[1]{\colorbox{yellow!35}{\texttt{\footnotesize #1}}}
\newcommand{\cardcheck}{\textsc{CardCheck}}
\newcommand{\kh}{\kappa_{\mathrm{HH}}}
\newcommand{\kv}{\kappa_{\mathrm{VH}}}
\newcommand{\dk}{\Delta\kappa}

\begin{document}

\title{The Delegation Map: Which Design-Card Constraints Can Be Checked
Automatically, and Which Only Look Like They Can}

\author{\IEEEauthorblockN{Anonymous Author(s)}
\IEEEauthorblockA{Submitted to ICSE 2027 --- double-blind review}}

\maketitle

\begin{abstract}
Coding agents increasingly produce repository-scale patches that must respect a
project's \emph{design cards}: lightweight, natural-language constraints on
architecture, API usage, naming, and design intent that no test enforces.
Checking every patch against every card by hand does not scale, and automated
verifiers are commonly summarized by a single aggregate accuracy against human
labels. We show that this summary is doubly broken: it hides that constraint
categories differ sharply in \emph{operationalizability}, and it treats human
review as a fixed oracle when human agreement itself varies by category. We
contribute a \emph{delegation map}: a per-category partition of design-card
constraints into \emph{automatable}, \emph{human-required},
\emph{reference-unstable}, and \emph{inconclusive}, built from a three-way-aligned
corpus of \ph{N} design cards, \ph{M} coding-agent patches across
\ph{R} repositories and \ph{A} agents, and independent human adjudications. Our
estimand is the per-category gap $\dk$ between the measured human--human
agreement ceiling (Fleiss'~$\kappa$) and verifier--human agreement
(Cohen's~$\kappa$); reference-unstable categories --- decided from human
agreement alone, before any verifier output is examined --- report the ceiling
only, since precision and recall against a reference annotators cannot agree on
are undefined by construction. Using a frozen hybrid verifier (\cardcheck) purely
as a measurement instrument, we find that \ph{K\_auto} structural categories
reach the human band (precision $\geq$ \ph{P*}, recall $\geq$ \ph{R*}, coverage
$\geq$ \ph{C*}), while intent-level categories sit \ph{DK} $\kappa$ below the
ceiling and \ph{K\_unstable} categories are reference-unstable. The resulting
routing rule cuts mandatory human review by \ph{RR}\% on our corpus mix and is
released as a per-category table from which any team can recompute savings for
its own card distribution. All criteria, thresholds, and falsifiers were fixed
before data collection; corpus, protocol, and instrument are released.
\end{abstract}

\begin{IEEEkeywords}
coding agents, design constraints, architecture conformance, LLM-as-judge,
inter-rater agreement, pre-registration
\end{IEEEkeywords}

% =========================================================================
\section{Introduction}\label{sec:intro}

Human review of design-card compliance is not a fixed oracle: its reliability
varies by constraint type, and any claim that automation ``matches human
review'' is meaningless until that variation is measured. A design card states
what an agent-generated patch may and may not do --- architectural boundaries it
must not cross, APIs it must not call, naming conventions it must follow, or
higher-level intent such as ``keep business logic separate from presentation.''
For a constraint like \emph{``do not import package X from module Y,''} two
reviewers will almost always agree, and a tool can compute the same verdict from
the patch's import graph. For \emph{``do not mix business logic with
presentation,''} reviewers themselves disagree, because the card names an intent
it does not fully encode. These are not two points on one difficulty scale; they
are two different verification problems, and the second lacks a stable reference
against which \emph{any} checker --- human or automated --- can be scored. This
paper measures that boundary.

As coding agents move from local edits to repository-scale
patches~\cite{jimenez2024swebench,yang2024sweagent}, manually checking every
patch against every card stops scaling, and the natural response is an automated
verifier that predicts violation from the patch. The field then summarizes such
a verifier with one aggregate accuracy number. That aggregate hides the only
decision a practitioner actually faces: \emph{which constraint categories can be
safely delegated to automation, which still require a human, and which are so
under-specified that even human review cannot serve as ground truth?} A verifier
with strong average accuracy can be unsafe because it passes structural checks
and silently misses intent violations; a verifier that looks weak in aggregate
may already be reliable on the categories that dominate a team's review load.
The right unit of analysis is the \emph{constraint category}, and the right
reference is not human labels taken as truth but the \emph{measured human--human
agreement ceiling} for that category.

We produce a \emph{delegation map}: a category-level partition of design-card
constraints into \emph{automatable}, \emph{human-required}, and
\emph{reference-unstable}, plus an \emph{inconclusive} cell where the evidence
does not yet decide. The category scheme is not invented for this study: it is
derived from the rule-type structure of existing architecture-conformance
tooling (dependency and layering rules in the style of
import-linter~\cite{importlinter} and ArchUnit~\cite{archunit}; API-usage and
deprecation rules; naming and location conventions; and the intent-level
principles those tools explicitly do \emph{not} encode), so the taxonomy has
provenance independent of our results, and we pre-register a
granularity-sensitivity check on it (Sec.~\ref{sec:design:stability}). The map
is built from a \emph{three-way-aligned corpus}: category-labeled design cards,
coding-agent patches scoped to those cards, and independent human adjudications
over the same card--patch pairs. For each category we estimate the gap between
human--human agreement (Fleiss'~$\kappa$~\cite{fleiss1971}, the ceiling) and
verifier--human agreement (Cohen's~$\kappa$~\cite{cohen1960}), and classify a
category as automatable only when that gap's confidence interval sits within a
pre-registered band of the ceiling. Categories classified
\emph{reference-unstable} report the human ceiling \emph{only}: their precision
and recall against human labels are left \emph{undefined by construction},
because a reference the annotators themselves cannot agree on cannot ground a
P/R number. Reference-unstable status is decided solely from human--human
$\kappa$, \emph{before any verifier output on that category is examined}, so the
exclusion cannot be steered by how the checker happens to perform.

The instrument that produces verifier labels, \cardcheck{}, applies
deterministic structural checks where a constraint's truth condition is
computable from the patch, and a calibrated, abstaining LLM judgment where it is
not. \cardcheck{} is not the contribution and is not tuned for leaderboard
superiority; it is the measurement device, held frozen so the map reflects the
constraints, not the checker. Because an instrument that abstains on every hard
item would earn a flattering $\kappa$ on the easy remainder, abstention is a
first-class outcome: every automatability verdict requires both the
$\kappa$-band \emph{and} a pre-registered coverage floor
(Sec.~\ref{sec:design:instrument}).

Our headline result is that operationalizability is sharply category-dependent.
On \ph{N} design cards and \ph{M} card--patch pairs across \ph{R} repositories
and \ph{A} coding agents, \ph{K\_auto} structural categories are automatable ---
verifier--human $\kappa$ within \ph{BAND} of the human ceiling, at precision
$\geq$ \ph{P*}, recall $\geq$ \ph{R*}, and coverage $\geq$ \ph{C*} --- while
intent-level constraints sit \ph{DK} $\kappa$ below the ceiling and are
classified human-required. The partition itself is stable: the automatable set
is invariant as the reference floor moves across \ph{KF}$\,\pm\,0.05$, with only
the human-required/reference-unstable boundary shifting by
$\leq$~\ph{SHIFT} categories, and no category changes verdict class under the
pre-registered merge/split granularity check. Routing human review only to
non-automatable categories cuts mandatory review load by \ph{RR}\% \emph{on our
corpus's category mix}; because we release the full per-category routing table
(category $\rightarrow$ verdict $\rightarrow$ $\dk$ $\rightarrow$ P/R where
defined), any team can recompute the reduction for its own card distribution.
The routed split survives a single-shot LLM-judge control budget-matched on
tokens. The gap on intent constraints is not weak engineering:
\ph{K\_unstable} categories are reference-unstable --- the humans themselves
fall below usable agreement --- so their failure mode is a card that
under-encodes its own truth condition, not a checker that can be improved.

\smallskip\noindent This paper contributes:
\begin{enumerate}[leftmargin=1.4em]
  \item \textbf{An estimand.} The per-category $\kappa$-gap between the human
  agreement ceiling and verifier--human agreement, which turns ``is automation
  good enough here?'' into a measurable, category-local quantity with a
  pre-registered decision band --- and which makes the reference-unstable case a
  first-class, reportable outcome rather than a source of noisy P/R.
  \item \textbf{A map, with a routing function.} A partition of
  \ph{N\_cat} design-card constraint categories into automatable /
  human-required / reference-unstable / inconclusive, shown stable under floor-
  and granularity-perturbation, delivered as a per-category routing table from
  which any team can recompute review savings for its own card mix.
  \item \textbf{A diagnostic (secondary).} For every human-required and
  reference-unstable card, a double-coded annotation of \emph{which}
  under-encoding it exhibits --- quantifier scope, implicit boundaries, or
  cross-file coupling --- with coder agreement reported, telling card authors
  what to encode to move a category toward automatable.
\end{enumerate}

% =========================================================================
\section{Background and Motivating Example}\label{sec:background}

\subsection{Design cards as governance objects}
Tests verify behavior; design cards encode architectural boundaries, security
assumptions, maintainability rules, and project-specific intent that tests do
not capture. An agent can produce a patch that compiles and passes every test
while violating the constraints that keep the codebase safe to change ---
the constraint classes that conformance research has targeted since reflexion
models~\cite{murphy1995reflexion,knodel2007comparison,passos2010static}, now
re-materialized as lightweight review gates for agent
workflows~\cite{sadowski2018lessons,tufano2021towards}. As teams delegate more
implementation to agents, the scarce resource stops being code generation and
becomes \emph{review}: deciding whether generated code respects rules no test
enforces.

\subsection{Running example: two cards, two problems}\label{sec:background:example}
Consider two cards from the same repository:
\begin{itemize}[leftmargin=1.2em]
  \item[$C_1$:] \emph{``Modules under \texttt{core/} must not import from
  \texttt{web/}.''}
  \item[$C_2$:] \emph{``Handlers must not mix business logic with presentation
  logic.''}
\end{itemize}
$C_1$'s truth condition is computable from the patch's import graph; two
reviewers and a ten-line script all return the same verdict. $C_2$ names an
intent: whether a given conditional in a handler is ``business logic'' depends
on a design rationale the card gestures at but does not encode. On such cards
reviewers disagree with each other --- and when they do, a verifier--human
disagreement is ambiguous between a verifier defect and a card defect. A single
uniform checker scored by a single aggregate number treats $C_1$ and $C_2$ as
one problem; a team that trusts the aggregate delegates exactly the checks it
should have kept.

\subsection{The reference must be measured, not assumed}
Human labels are the natural reference for compliance, but on intent-laden
cards, annotators disagree, and chance-corrected agreement is the standard
instrument for quantifying how much~\cite{artstein2008inter,mchugh2012kappa}.
Measuring the human--human ceiling per category resolves the ambiguity of a
verifier--human disagreement into three interpretable cases: automation reaches
the human band (\emph{automatable}); automation falls short of a usable band
(\emph{human-required}); or humans do not themselves reach a usable band
(\emph{reference-unstable}). The third case is the one existing evaluations
mishandle. When annotators cannot agree, precision and recall against their
labels are not merely noisy --- they are \emph{undefined}, because there is no
reference to be precise or recall-complete against. Reporting a P/R number
there manufactures false confidence. We make the reference-unstable verdict
visible machinery with two safeguards fixed in advance
(Sec.~\ref{sec:design:verdicts}).

% =========================================================================
\section{Study Design}\label{sec:design}

\subsection{Research questions}\label{sec:design:rqs}
\begin{itemize}[leftmargin=1.6em]
  \item[\textbf{RQ1}] (\emph{Operationalizability gap}) Do structural constraint
  categories reach the human-agreement ceiling while intent-level categories
  remain measurably below it?
  \item[\textbf{RQ2}] (\emph{Automatable quality}) On categories within the
  human band, does the frozen instrument meet absolute precision/recall/coverage
  floors and survive a token-budget-matched LLM-judge control?
  \item[\textbf{RQ3}] (\emph{Partition stability}) Is the delegation map
  invariant under perturbation of the reference floor and of the taxonomy's
  granularity?
  \item[\textbf{RQ4}] (\emph{Under-encoding diagnostic; secondary}) Which
  under-encoding types explain the categories that resist automation?
\end{itemize}

\subsection{Estimand and verdict rules}\label{sec:design:verdicts}
For category $c$, let $\kh(c)$ be Fleiss'~$\kappa$~\cite{fleiss1971} over
independent human labels and $\kv(c)$ be Cohen's~$\kappa$~\cite{cohen1960}
between the frozen verifier and the human majority label, computed over the
items the verifier did not abstain on. The estimand is
\begin{equation}
  \dk(c) \;=\; \kh(c) \;-\; \kv(c),
\end{equation}
with confidence intervals from a cluster bootstrap over cards (card--patch
pairs nest within cards; resampling pairs i.i.d.\ would understate variance).
Verdicts are assigned by pre-registered rules, in this order:

\begin{enumerate}[leftmargin=1.5em]
  \item \textbf{Reference-unstable (verifier-blind).} If
  $\kh(c) < $ \ph{KF} (anchored to the Landis--Koch ``moderate''
  boundary, $\kappa=0.40$~\cite{landis1977}), category $c$ is
  reference-unstable. This is decided from human labels alone, before any
  verifier output on $c$ is examined. The category reports $\kh(c)$ only;
  its P/R are undefined and omitted from every automatability claim. A
  transparency appendix reports what the instrument's P/R \emph{would have
  been}, labeled strictly as undefined-reference diagnostics that ground no
  claim and enter no comparison.
  \item \textbf{Automatable.} $c$ is automatable iff the CI upper bound of
  $\dk(c)$ is $\leq$ \ph{BAND}, \emph{and} precision $\geq$ \ph{P*} and recall
  $\geq$ \ph{R*} against the human majority label, \emph{and} the instrument's
  coverage (fraction of items not abstained) is $\geq$ \ph{C*}.
  \item \textbf{Human-required.} $c$ is human-required iff the CI lower bound of
  $\dk(c)$ exceeds \ph{BAND}, or any absolute floor in rule~2 fails while
  $\kh(c)$ is usable.
  \item \textbf{Inconclusive.} Otherwise --- the CI straddles the band or the
  cell is too thin to decide. Inconclusive is a reported verdict, not a
  category silently dropped.
\end{enumerate}

Because $\kappa$ is prevalence-sensitive~\cite{mchugh2012kappa}, each category
additionally reports raw percent agreement and prevalence-adjusted agreement
(PABAK) as pre-registered robustness columns; $\kappa$ remains primary and no
verdict is revised on the robustness columns alone.

\wj{Standpoint for the verdict machinery: the verifier-blind ordering in rule~1
is the paper's integrity core. We expect a hostile reviewer to read
``P/R undefined by construction'' as ``the authors excluded their worst
categories.'' Our answer is procedural, not rhetorical: the exclusion is
computable from human labels alone, its threshold is anchored to a
50-year-old convention rather than chosen by us, and the would-have-been
numbers are printed anyway in the appendix. If the appendix numbers turn out
\emph{good} on unstable categories, that is fine too --- it strengthens the
claim that the exclusion is principled, since we excluded categories where the
checker looked strong.}

\subsection{Constraint taxonomy with external provenance}\label{sec:design:taxonomy}
The category scheme is derived from the rule-type structure of existing
conformance tooling, not invented post hoc: dependency/layering rules and
API-usage/deprecation rules follow the rule vocabularies of
import-linter~\cite{importlinter}, ArchUnit~\cite{archunit}, and the static
conformance literature~\cite{knodel2007comparison,passos2010static};
naming/location conventions follow linter rule classes; the intent-level
remainder (separation of concerns, coupling discipline, abstraction
boundaries) is precisely the class those tools document as out of scope.
The full mapping from tool rule types to our \ph{N\_cat} categories is given in
the artifact.
% TODO-verify: if the Dong et al. constraint taxonomy checks out, cite it here
% as convergent provenance; do NOT cite before verification.

\subsection{Three-way-aligned corpus}\label{sec:design:corpus}
The corpus aligns three layers over the same card--patch pairs.
\emph{(i) Cards:} \ph{N} design cards drawn from \ph{R} repositories, each
labeled with a taxonomy category by two coders (disagreements adjudicated;
coder agreement reported).
\emph{(ii) Patches:} \ph{M} card--patch pairs produced by \ph{A} coding agents
given repository-scoped tasks that touch card-relevant code. Violations arise
two ways, both pre-registered: \emph{natural} violations (the agent is never
shown the card) and \emph{induced} variants (the task is chosen so that the
straightforward implementation conflicts with the card, again without showing
the card). The natural violation rate \ph{NVR} and the induced share are
reported; no patch is hand-edited.
\emph{(iii) Human labels:} each pair receives $\geq 3$ independent binary
violation labels with confidence and a one-sentence rationale, from a pinned
annotator pool \ph{POOL} (composition fixed before the pilot;
Sec.~\ref{sec:threats}). Annotators see the card, the patch, and surrounding
context, but never the verifier's output.

\subsection{The frozen instrument: \cardcheck{}}\label{sec:design:instrument}
\cardcheck{} maps a (card, patch) pair to \{\emph{violation},
\emph{no-violation}, \emph{abstain}\}. Structurally checkable cards route to
deterministic checkers over the patch's import graph, call sites, AST patterns,
and file layout; remaining cards route to a calibrated LLM judgment with an
explicit abstain option under low confidence. Three design rules make it an
instrument rather than a contribution:
\emph{(a) frozen:} all prompts, thresholds, and routing rules are fixed during
an exploratory tuning phase on a disjoint tuning split and never changed during
the confirmatory run;
\emph{(b) abstention-honest:} abstentions are excluded from $\kv$ but counted
against the coverage floor \ph{C*}, so selective abstention cannot buy
agreement;
\emph{(c) budget-accounted:} every call's token count is logged, enabling the
budget-matched control.

\wj{Standpoint on the instrument: we deliberately do not claim
\cardcheck{} is better than anything. The hybrid design exists to make the map's
verdicts attributable to the \emph{constraints} rather than to instrument
weakness --- deterministic checks remove LLM noise where structure suffices,
and abstention with a coverage floor removes the ``grade only the easy items''
gaming path. If a reviewer asks ``why not a stronger verifier?'', the answer
is: a stronger verifier can only shrink $\dk$ on categories with a usable
reference; it cannot create a reference where humans have none, which is the
boundary the paper measures.}

\subsection{Baselines and budget-matched control}\label{sec:design:baselines}
Three comparison points, all evaluated on the identical pairs:
\emph{floor} --- a naive keyword/regex checker compiled from each card's
surface text;
\emph{ceiling} --- the human majority label (by construction $\kv \leq \kh$ in
expectation);
\emph{control} --- a single-shot LLM judge given the card and patch, matched to
\cardcheck{} on \textbf{tokens} (primary unit; matched LLM-call count reported
as a robustness column). The control answers the cheapest rival explanation:
that any routed hybrid ``wins'' only by spending more.
McNemar's test on paired verdicts compares \cardcheck{} to the control within
each usable category, with Holm correction across the category family.

\subsection{Partition stability checks}\label{sec:design:stability}
Two perturbations are pre-registered.
\emph{Floor sensitivity:} the entire partition is recomputed at
\ph{KF}$-0.05$, \ph{KF}, and \ph{KF}$+0.05$. Claim form: the automatable set is
invariant across the band; only the human-required$\,\leftrightarrow\,$%
reference-unstable boundary may shift, by $\leq$ \ph{SHIFT} categories.
\emph{Granularity:} the two most similar categories are merged and the largest
category is split (rules fixed in the artifact), the classification is re-run,
and every verdict-class migration is reported. If verdicts are
granularity-stable, the map is real at this resolution; if not, the migration
count bounds how fine a routing rule can be.

\subsection{Pre-registered criteria and falsifiers}\label{sec:design:prereg}
The central claim holds \textbf{iff}:
(1)~$\geq$ \ph{K\_struct} structural categories have $\dk$ CI upper bound
$\leq$ \ph{BAND};
(2)~the intent-level categories have $\dk$ CI lower bound $>$ \ph{BAND};
(3)~automatable categories reach precision $\geq$ \ph{P*}, recall $\geq$
\ph{R*}, coverage $\geq$ \ph{C*};
(4)~the automatable/human-required split survives the token-budget-matched
control;
(5)~the automatable set is invariant across the $\kappa$-floor sensitivity
band.
The claim is \textbf{falsified} if
(a)~all categories show similar $\dk$;
(b)~intent categories also reach the human band;
(c)~$\kh$ is below the usable floor in so many categories that no stable
partition exists; or
(d)~the automatable set is not invariant across the sensitivity band.
Cases (c) and (d) are findings, not failed experiments: (c) says design cards
under-encode as a format; (d) bounds the granularity at which any routing rule
can exist~\cite{nosek2018preregistration}.

\wj{Standpoint on the criteria: (1)+(2) is the two-sided form of H1 --- we are
not merely predicting ``structural is easier,'' which is folklore, but that the
gap is large enough to clear a pre-registered band on both sides, with CIs. The
band \ph{BAND} is the one threshold still deliberately unbound: it will be set
from the pilot's human-$\kappa$ dispersion (derivation to be written into
paper.md before the confirmatory run), because hardcoding it now, before we
know human-ceiling variance, would be pseudo-rigor.}

% =========================================================================
\section{Experiment Setup}\label{sec:setup}

\subsection{Repositories, agents, and scale}
We select \ph{R} repositories spanning \ph{LANGS} languages and at least three
architecture styles (layered, hexagonal/ports-adapters, plugin-based), each with
enough conformance-relevant structure to author cards against; repository
selection criteria are fixed in the artifact. Patches come from \ph{A} coding
agents \ph{AGENT\_LIST}, each run under its default configuration with
repository-scoped tasks of \ph{TASKS\_PER\_REPO} tasks per repository. The final
corpus comprises \ph{N} cards ($\geq$ \ph{MIN\_PER\_CAT} per category) and
\ph{M} card--patch pairs.

\subsection{Annotation protocol}
Each pair is labeled independently by three annotators from the pinned pool
\ph{POOL}; order is randomized, cards from the same repository are never
consecutive for the same annotator, and no annotator labels a pair from a
repository they contributed code to. Fleiss'~$\kappa$ is computed from the raw
independent labels; the majority label serves as the P/R reference only in
categories that pass the usable-reference floor. For RQ4, every human-required
and reference-unstable card is double-coded against the three under-encoding
types (quantifier scope, implicit boundaries, cross-file coupling); if coder
agreement on this secondary annotation falls below \ph{KF}, the diagnostic is
demoted to qualitative discussion by pre-registered rule.

\subsection{Tuning/confirmatory split and pilot}
All \cardcheck{} tuning happens on a disjoint tuning split
(\ph{TUNE\_FRAC} of cards, stratified by category) and is declared exploratory.
The confirmatory run executes once, on the frozen instrument, on the remaining
cards. Before the full corpus, a two-category pilot (API constraints vs.\
design-principle constraints; 10 cards, 40 pairs, 3 annotators) must
mechanically produce per-category $\kappa$/P/R, run the floor-sensitivity
check, and reach usable human $\kappa$ in at least one category --- the
pre-registered gate for scaling.

\subsection{Statistical analysis}
CIs: cluster bootstrap over cards, \ph{B} resamples. Verifier vs.\ control:
McNemar within category, Holm-corrected across categories. All tests, the
$\alpha$ level \ph{ALPHA}, and the handling of thin cells ($<$
\ph{MIN\_CELL} pairs $\Rightarrow$ inconclusive by rule) are fixed here, before
data.

% =========================================================================
\section{Results}\label{sec:results}

\noindent\emph{Every number below is a pre-registered placeholder (\ph{~}); the
hypotheses and their standpoints are stated now, before data, so that any
outcome maps to committed wording.}

\subsection{RQ1: The operationalizability gap}
\textbf{H1} (criteria 1--2): at least \ph{K\_struct} structural categories have
$\dk$ CI upper bound $\leq$ \ph{BAND}, while intent categories have $\dk$ CI
lower bound $>$ \ph{BAND}.

Table~\ref{tab:map} reports, per category: $\kh$, $\kv$, $\dk$ with CI, percent
agreement, PABAK, coverage, and verdict.
Structural categories show $\dk =$ \ph{DK\_struct\_range}; intent categories
show $\dk =$ \ph{DK\_intent\_range}; \ph{K\_unstable} categories fall below the
usable floor and report $\kh$ only.

\wj{H1 standpoint: this is the load-bearing hypothesis and it is genuinely
two-sided. If intent categories also reach the band (falsifier b), the paper
inverts into ``design-card checking is more automatable than believed'' --- a
publishable surprise with committed wording. If everything shows similar
$\dk$ (falsifier a), the category axis is not where operationalizability
lives, and we stop. The expected outcome, from the pilot's shape, is a clean
split; the risk we rate highest is not the split failing but too many
categories landing reference-unstable, which flips us into case (c):
the format-level under-encoding finding. We commit to reporting whichever
fires.}

\begin{table}[t]
\caption{The delegation map (per-category; placeholder values pending the
confirmatory run). RU = reference-unstable: $\kh$ below \ph{KF}; P/R undefined
by construction, reported only in the diagnostics appendix.}
\label{tab:map}
\centering
\footnotesize
\begin{tabular}{@{}lcccccc@{}}
\toprule
Category & $\kh$ & $\kv$ & $\dk$ [CI] & Cov. & P / R & Verdict \\
\midrule
Dep./layering   & \ph{x} & \ph{x} & \ph{x} & \ph{x} & \ph{x}/\ph{x} & \ph{Auto} \\
API usage       & \ph{x} & \ph{x} & \ph{x} & \ph{x} & \ph{x}/\ph{x} & \ph{Auto} \\
Naming/location & \ph{x} & \ph{x} & \ph{x} & \ph{x} & \ph{x}/\ph{x} & \ph{Auto} \\
Security invar. & \ph{x} & \ph{x} & \ph{x} & \ph{x} & \ph{x}/\ph{x} & \ph{?} \\
Sep.\ of concerns & \ph{x} & \ph{x} & \ph{x} & \ph{x} & --- & \ph{HR/RU} \\
Coupling/abstr. & \ph{x} & \ph{x} & \ph{x} & \ph{x} & --- & \ph{HR/RU} \\
\bottomrule
\end{tabular}
\end{table}

\subsection{RQ2: Quality on automatable categories}
\textbf{H2} (criteria 3--4): automatable categories meet precision $\geq$
\ph{P*}, recall $\geq$ \ph{R*}, coverage $\geq$ \ph{C*}, improve F1 over the
regex floor by \ph{F1\_gain}, and the split survives the token-matched control
(McNemar, Holm-corrected: \ph{pvals}).

\wj{H2 standpoint: the control is the criterion most likely to bite. If the
single-shot judge matches \cardcheck{} within budget on structural categories,
the honest conclusion is that deterministic checking buys nothing beyond a
prompt --- which weakens the instrument story but leaves the map (contribution
1--2) intact, since the map is about the reference, not the checker. We
pre-commit to that wording rather than defending the hybrid.}

\subsection{RQ3: Partition stability}
\textbf{H3} (criterion 5): the automatable set is invariant across
\ph{KF}$\pm0.05$ (boundary shift $\leq$ \ph{SHIFT} categories), and no category
changes verdict class under the merge/split granularity check
(migrations: \ph{MIG}).

\wj{H3 standpoint: instability here is not fatal --- falsifier (d) converts it
into a granularity bound (``routing rules coarser than $k$ categories are the
finest supportable''). What we cannot survive is discovering instability only
at rebuttal; hence the pilot runs this check mechanically before scaling.}

\subsection{RQ4: Under-encoding diagnostic (secondary)}
\textbf{H4}: among human-required and reference-unstable cards, the
distribution over \{quantifier scope, implicit boundaries, cross-file
coupling\} is \ph{DIST}, with double-coder agreement \ph{KD}.

\wj{H4 standpoint: strictly secondary and pre-registered to demote itself
(Sec.~\ref{sec:setup}) if coder agreement on the under-encoding annotation is
itself unusable --- it would be self-refuting to hand-wave the diagnostic in a
paper about unstable references. Its value if it survives: it converts the
map's negative cells into an actionable card-authoring checklist.}

\subsection{Routing yield}
Under the map, routing human review only to non-automatable categories cuts
mandatory review load by \ph{RR}\% at retained recall $\geq$ \ph{RC}\% on
automatable categories, \emph{for our corpus's category mix}
(\ph{MIX}). The released routing table makes the savings a function of the
reader's card distribution rather than a corpus-specific percentage.

% =========================================================================
\section{Discussion}\label{sec:discussion}
\textbf{From leaderboard to routing rule.} The practical reading of
Table~\ref{tab:map} is procedural: send dependency, API, and naming cards to the
tool; keep intent cards with humans; and treat reference-unstable cards as
\emph{card-authoring} bugs --- the fix is encoding the truth condition
(RQ4 names which one), not buying a better checker.
\textbf{What the map does not claim.} Verdicts are claims about this reviewer
population, this taxonomy resolution (bounded by RQ3), and this instrument
generation; the released table and protocol let any team re-derive all three
locally.
\textbf{Relation to agent-evaluation gaming.} The abstention/coverage
machinery is a small instance of a general principle: any evaluator with a
discretionary opt-out must be scored jointly on agreement and coverage, or the
opt-out becomes the gaming surface.

% =========================================================================
\section{Threats to Validity}\label{sec:threats}
\textbf{Construct.} $\kappa$ is prevalence-sensitive; we report percent
agreement and PABAK alongside, pre-registered as robustness-only. The majority
label is a reference, not truth; that is precisely why verdicts are phrased as
gaps to a measured ceiling rather than accuracy against ground truth.
\textbf{Internal.} Instrument freeze is enforced by a disjoint tuning split and
released hashes; annotators never see verifier output; reference-unstable
status is computed verifier-blind.
\textbf{External.} The ceiling is the ceiling \emph{for our pinned pool}
\ph{POOL}; a category unstable for external reviewers might be stable for a
project's own maintainers. The map generalizes only as far as \ph{R}
repositories, \ph{A} agents, and the taxonomy's resolution; RQ3 bounds the
latter, and the routing table lets teams re-derive verdicts with their own
reviewers.
\textbf{Conclusion.} Per-category tests are Holm-corrected; thin cells
($<$ \ph{MIN\_CELL} pairs) are inconclusive by rule, never silently dropped.

% =========================================================================
\section{Related Work}\label{sec:related}
\textbf{Architecture conformance checking.} Reflexion
models~\cite{murphy1995reflexion} and their successors check code against
architectural rules, with comparative studies of static
approaches~\cite{knodel2007comparison,passos2010static} and industrial linting
at scale~\cite{sadowski2018lessons}; modern rule engines (ArchUnit, import-%
linter)~\cite{archunit,importlinter} mechanize exactly the structural rule
types our taxonomy inherits. This line assumes the rule's truth condition is
already encoded; it does not ask which natural-language constraints \emph{can}
be so encoded, which is our estimand.
\textbf{LLM-as-judge and its reliability.} LLM judges are widely used to score
generation~\cite{zheng2023judging}, and agent benchmarks increasingly rely on
automated graders~\cite{jimenez2024swebench,yang2024sweagent}. Reliability
concerns are typically addressed by scaling or ensembling judges; our point is
prior to that: for correctness-style tasks the human reference itself must be
measured per category, since no ensemble fixes a reference that does not exist.
\textbf{Agreement measurement.} We import chance-corrected agreement and its
known pathologies from content
analysis~\cite{cohen1960,fleiss1971,landis1977,artstein2008inter,mchugh2012kappa};
the delegation-map construction --- verdicts defined by a gap to a measured
ceiling, with an undefined-reference cell --- is, to our knowledge, new in
SE evaluation of automated review.
% TODO-verify: add 2024--2026 LLM-code-review-agreement refs here after
% checking each against its venue page; none are cited on memory.
\textbf{Idea-shape and pre-registration.} Our confirmatory design follows
pre-registration practice~\cite{nosek2018preregistration}; the framing
discipline (mechanism-level moves over bridge/synthesis) follows the measured
gap between human and LLM research ideas~\cite{chen2026measuring}.

% =========================================================================
\section{Conclusion}\label{sec:conclusion}
We reframed automated design-card checking from ``can we build a better
verifier?'' to ``for which kinds of constraints is verification delegable at
all?'' --- and made the second question measurable via a per-category
$\kappa$-gap to a human ceiling that is measured rather than assumed, with an
undefined-reference cell for the cards on which no checker, human or machine,
can currently be scored. The resulting delegation map, its stability checks,
and its per-category routing table are released so that both the verdicts and
the savings can be recomputed by any team, for any card mix, with any reviewer
pool. \ph{One-sentence headline number, bound to H1, inserted after the
confirmatory run.}

\balance
\bibliographystyle{IEEEtran}
\begin{thebibliography}{20}
\bibitem{cohen1960} J. Cohen, ``A coefficient of agreement for nominal
scales,'' \emph{Educ. Psychol. Meas.}, 20(1):37--46, 1960.
\bibitem{fleiss1971} J. L. Fleiss, ``Measuring nominal scale agreement among
many raters,'' \emph{Psychol. Bull.}, 76(5):378--382, 1971.
\bibitem{landis1977} J. R. Landis and G. G. Koch, ``The measurement of observer
agreement for categorical data,'' \emph{Biometrics}, 33(1):159--174, 1977.
\bibitem{artstein2008inter} R. Artstein and M. Poesio, ``Inter-coder agreement
for computational linguistics,'' \emph{Comput. Linguist.}, 34(4):555--596,
2008.
\bibitem{mchugh2012kappa} M. L. McHugh, ``Interrater reliability: the kappa
statistic,'' \emph{Biochemia Medica}, 22(3):276--282, 2012.
\bibitem{murphy1995reflexion} G. C. Murphy, D. Notkin, and K. Sullivan,
``Software reflexion models: bridging the gap between source and high-level
models,'' in \emph{Proc. FSE}, 1995.
\bibitem{knodel2007comparison} J. Knodel and D. Popescu, ``A comparison of
static architecture compliance checking approaches,'' in \emph{Proc. WICSA},
2007.
\bibitem{passos2010static} L. Passos, R. Terra, M. T. Valente, R. Diniz, and
N. das Chagas Mendon\c{c}a, ``Static architecture-conformance checking: an
illustrative overview,'' \emph{IEEE Software}, 27(5):82--89, 2010.
\bibitem{sadowski2018lessons} C. Sadowski, E. Aftandilian, A. Eagle,
L. Miller-Cushon, and C. Jaspan, ``Lessons from building static analysis tools
at Google,'' \emph{Commun. ACM}, 61(4):58--66, 2018.
\bibitem{tufano2021towards} R. Tufano, L. Pascarella, M. Tufano, D. Poshyvanyk,
and G. Bavota, ``Towards automating code review activities,'' in \emph{Proc.
ICSE}, 2021.
\bibitem{zheng2023judging} L. Zheng \emph{et al.}, ``Judging LLM-as-a-judge
with MT-Bench and Chatbot Arena,'' in \emph{Proc. NeurIPS Datasets and
Benchmarks}, 2023.
\bibitem{jimenez2024swebench} C. E. Jimenez \emph{et al.}, ``SWE-bench: can
language models resolve real-world GitHub issues?'' in \emph{Proc. ICLR},
2024.
\bibitem{yang2024sweagent} J. Yang \emph{et al.}, ``SWE-agent: agent--computer
interfaces enable automated software engineering,'' in \emph{Proc. NeurIPS},
2024.
\bibitem{nosek2018preregistration} B. A. Nosek, C. R. Ebersole, A. C. DeHaven,
and D. T. Mellor, ``The preregistration revolution,'' \emph{PNAS},
115(11):2600--2606, 2018.
\bibitem{chen2026measuring} Chen, Zhao, and Cohan, ``Measuring the gap between
human and LLM research ideas,'' arXiv:2607.01233, 2026.
\bibitem{archunit} ``ArchUnit: unit test your Java architecture,''
\url{https://www.archunit.org}.
\bibitem{importlinter} ``Import Linter: enforce contracts between Python
packages,'' \url{https://github.com/seddonym/import-linter}.
\end{thebibliography}

\end{document}
